\section{Conclusion, Limitations, and Future Directions}
\label{sec:conclusion_future}

This technical report presents a first-stage practice for full-parameter post-training of the DeepSeek-V4 model family on Ascend NPU SuperPOD. We refer to the overall framework as \textbf{\modelname}, which couples full-stack Ascend SuperPOD training optimization with solver-grounded CPT--SFT specialization for OR. 

At the system level, we demonstrate one of the first full-stack optimization efforts for trillion-parameter MoE LLM training on a non-GPU computing platform. The proposed end-to-end optimization strategy spans system-level training orchestration and bottleneck-level kernel optimization, achieving up to 34.22\% MFU. These results validate both the training efficiency and practical feasibility of large-scale post-training on Ascend NPU SuperPOD.

For post-training tasks, we establish an OR-oriented CPT--SFT workflow on DeepSeek-V4-Flash. Under identical SFT settings, direct SFT improves B4O-Feasible from 60.47\% to 65.93\% and B4O-ORGEval from 34.26\% to 48.73\%. With CPT initialization, the corresponding scores further increase to 71.22\% and 59.39\%, respectively. The full benchmark results show consistent trends across NL4OPT and OptiBench. These results suggest that direct SFT provides supervised task alignment, while CPT initialization introduces additional OR-domain modeling priors that improve both solver-facing feasibility and structural equivalence. The most significant improvement is observed on B4O-ORGEval, highlighting the value of domain-adaptive pre-training for complex OR modeling.



The SFT experiments further highlight the importance of data quality. Scaling self-distilled data from 10K to 50K does not lead to monotonic improvements and may even degrade performance on natural-language modeling benchmarks. In contrast, contract-aware cleaning and concise CoT-oriented enhancement improve the 10K data more reliably, especially on NL4OPT, OptiBench, and ORGEval. This indicates that OR-domain SFT benefits more from targeted repair of modeling errors, strict output-contract enforcement, compact reasoning checklists, and solver-grounded verification than from simply increasing the volume of synthetic training data.
Admittedly, the current post-training experiments are conducted on DeepSeek-V4-Flash. In future work, we will extend the same training pipeline to DeepSeek-V4-Pro and release the corresponding training results to further validate the scalability of the proposed infrastructure and post-training workflow.


\paragraph{Reproducibility and artifact release.}
To support full reproducibility and facilitate follow-up research, we  release the key artifacts underlying the proposed optimization and training workflow, including prompt templates, format-contract specifications, curriculum scheduling configurations, and monitoring callbacks. The AI-assisted data and evaluation tools used in our pipeline, including Cleaner, Reviewer, and Diagnostic Resolver, are implemented as configurable prompt- and rule-driven modules. Their templates and configuration files also are released, allowing future users to inspect, reproduce, adapt, and extend the data-cleaning, sample-review, diagnostic-resolution, and monitoring procedures introduced in this work.

In addition, we plan to separately release the optimized AscendC kernels developed through our kernel-optimization pipeline. This release will provide concrete kernel-level implementations and optimization examples, enabling future work to reproduce the reported operator-level improvements and further adapt the optimized kernels to related Ascend-based training workloads.

\paragraph{Future directions.}
Building on the proposed high-performance training infrastructure for Ascend SuperPOD, future work will further improve the efficiency of end-to-end full-parameter training while extending the post-training workflow toward agentic reinforcement learning~(AgenticRL). In particular, we will investigate operations-research-based mathematical modeling for training parallelism optimization, aiming to automate the exploration of multi-dimensional parallel configurations under memory, communication, and orchestration constraints. In parallel, we will continue scaling OR-oriented continued pre-training~(CPT) and supervised fine-tuning~(SFT) to broader solver-verified corpora, extend evaluation to more diverse OR problem families and general-capability benchmarks, and strengthen structural verification for nonlinear expressions, ratio constraints, quadratic terms, and Gurobi-specific modeling patterns.

Overall, this report demonstrates Ascend NPU SuperPOD can
support stable full-parameter post-training of trillion-scale models with MoE and hybrid attention architectures. 
Furthermore, the proposed CPT-SFT workflow can produce measurable gains in
OR-oriented mathematical modeling. The key finding is that effective vertical
domain adaptation depends on the alignment among domain knowledge acquisition,
solver-grounded data construction, supervised data quality, system stability,
and task-specific evaluation.

% \newpage
